\uuid{QlYZ}
\exo7id{4041}
\titre{exo7 4041}
\auteur{quercia}
\organisation{exo7}
\datecreate{2010-03-11}
\isIndication{false}
\isCorrection{true}
\chapitre{Développement limité}
\sousChapitre{Développements limités implicites}
\cpge{mpsi-mp2i.nombres-reels-suites.c, pcsi.nombres-reels-suites.b, ptsi.nombres-reels-suites.b, mpsi-mp2i.analyse-asymptotique.c, pcsi.analyse-asymptotique.c, ptsi.analyse-asymptotique.c, mpsi-mp2i.analyse-asymptotique.d, pcsi.analyse-asymptotique.d, bcpst1.suites-reelles, bcpst1.developpements-limites.a}
\module{Analyse}
\niveau{L1}
\difficulte{}

\contenu{
\question{
Soit $u_n$ une suite réelle telle que pour tout $n$ on ait 
$u_n^5+nu_n-1=0$. Trouver un développement asymptotique à deux termes de $u_n$.
}
\reponse{
$u_n = \frac1n - \frac1{n^6} +  o\Bigl(\frac1{n^6}\Bigr)$.
}
}
